\documentclass{article}
\usepackage{amssymb,amsmath}
\newtheorem{definition}{Definition}
\begin{document}
\title{Specification}
\author{}
\date{}
\maketitle
This document provides the definitions of multilattices, residuated lattices and residuated multilattices, and the respective notions of homomorphisms and
derivations, that are implemented by the program. As detailed in Sect.~\ref{rela0}, we consider residuated lattices to be bounded and commutative (with the
top element being the neutral element w.r.t. multiplication). The residuated multilattices in Sect.~\ref{remu0} are adapted from~\cite{CCG14}, but unlike there,
we consider them to be bounded (to match the previous definition of residuated lattices). The notion of multilattices in Sect.~\ref{mult0} is also taken
from~\cite{CCG14}. The notion of derivations on multilattices in Sect.~\ref{mult0} is based on unpublished work (as of December 2024)
by Darline Laure Keubeng Yemene. The notions of $(f,g)$-multiplicative and $(f,g)$-implicative derivations on residuated lattices,
in Sect.~\ref{rela0}, are obtained from~\cite{YFAL20}. The respective notion of $(f,g)$-multiplicative derivations on residuated multilattices,
in Sect.~\ref{remu0}, has been introduced in~\cite{YFL22}, there simply called $(f,g)$-derivation; whereas the notion of $(f,g)$-implicative derivations
on residuated multilattices, in Sect.~\ref{remu0}, should be considered tentative, having not yet been sufficiently explored.

\section{Multilattices}
\label{mult0}
Let $(P,\leq)$ be a poset. For any $X \subseteq P$, the elements of
\begin{align}
X^{\uparrow} &:= \{p \in P \mid p \geq x \text{ for all } x \in X\} \;,\\
X^{\downarrow} &:= \{p \in P \mid p \leq x \text{ for all } x \in X\}
\end{align}
are the \emph{upper bounds} and \emph{lower bounds} of $X$, respectively. For any $x,y \in P$, a minimal element of $\{x,y\}^{\uparrow}$
is a \emph{multisupremum} of $x$ and $y$, whereas a maximal element of $\{x,y\}^{\downarrow}$ is a \emph{multiinfimum} of $x$ and $y$. We use $x \sqcup y$ and
$x \sqcap y$ to denote the set of all multisuprema and multiinfima of $x$ and $y$, respectively.

\begin{definition}[Ordered Multilattice]
A poset $(M,\leq)$ is an \emph{ordered multilattice} if, for all $x,y \in M$, every element of $\{x,y\}^{\uparrow}$ lies above an element of $x \sqcup y$,
and every element of $\{x,y\}^{\downarrow}$ lies below an element of $x \sqcap y$.
\end{definition}
It should be noted that every finite poset $(M,\leq)$ is an ordered multilattice.

\begin{definition}[Multilattice]
The algebra $(M,\sqcup,\sqcap)$ associated with an ordered multilattice $(M,\leq)$ is a \emph{multilattice}.
\end{definition}

\begin{definition}[Bounded Multilattice]
The algebra $(M,\sqcup,\sqcap,0,1)$ associated with an ordered multilattice $(M,\leq)$ with least element $0$ and greatest element $1$
is a \emph{bounded multilattice}.
\end{definition}

\begin{definition}[Homomorphism]
A \emph{homomorphism} $f:(M_1,\sqcup,\sqcap) \rightarrow (M_2,\sqcup,\sqcap)$ of multilattices is a function $f:M_1 \rightarrow M_2$ which satisfies
$f(x \sqcup y) \subseteq f(x) \sqcup f(y)$ and $f(x \sqcap y) \subseteq f(x) \sqcap f(y)$ for all $x,y \in M_1$.
\end{definition}

\begin{definition}[Derivation]
Let $(M,\sqcup,\sqcap)$ be a multilattice. Moreover, let $f$ and $g$ be endomorphisms of $(M,\sqcup,\sqcap)$.
An \emph{$(f,g)$-derivation} of $(M,\sqcup,\sqcap)$ is a function $d:M \rightarrow M$ which satisfies
\begin{align}
d(x \sqcap y) \subseteq (d(x) \sqcap f(y)) \sqcup (g(x) \sqcap d(y))
\end{align}
for all $x,y \in M$.
\end{definition}

\section{Residuated Lattices}
\label{rela0}
\begin{definition}[Residuated Lattice~\cite{Pic21}]
A \emph{residuated lattice} is an algebra $(L,\vee,\wedge,\odot,\rightarrow,0,1)$ such that $(L,\vee,\wedge,0,1)$
is a bounded lattice, $(L,\odot,1)$ is a commutative monoid, and $\rightarrow$ is a binary operation satisfying
\begin{align}
x \odot z \leq y \Leftrightarrow x \leq z \rightarrow y
\end{align}
for all $x,y,z \in L$.
\end{definition}

\begin{definition}[Homomorphism]
A \emph{homomorphism} $f:(L_1,\vee,\wedge,\odot,\rightarrow,0,1) \rightarrow (L_2,\vee,\wedge,\odot,\rightarrow,0,1)$ of residuated lattices
is a function $f:L_1 \rightarrow L_2$ which satisfies
\begin{align}
f(x \vee y) &= f(x) \vee f(y) \\
f(x \wedge y) &= f(x) \wedge f(y) \\
f(x \odot y) &= f(x) \odot f(y) \\
f(x \rightarrow y) &= f(x) \rightarrow f(y) \\
f(0) &= 0 \\
f(1) &= 1
\end{align}
for all $x,y \in L_1$.
\end{definition}

\begin{definition}[Multiplicative Derivation~\cite{YFAL20}]
Let $\mathcal{L} := (L,\vee,\wedge,\odot,\rightarrow,0,1)$ be a residuated lattice. Moreover, let $f$ and $g$ be endomorphisms of $\mathcal{L}$.
An \emph{$(f,g)$-multiplicative derivation} of $\mathcal{L}$ is a function $d:L \rightarrow L$ which satisfies
\begin{align}
d(x \odot y) = (d(x) \odot f(y)) \vee (g(x) \odot d(y))
\end{align}
for all $x,y \in L$.
\end{definition}

\begin{definition}[Implicative Derivation~\cite{YFAL20}]
Let $\mathcal{L} := (L,\vee,\wedge,\odot,\rightarrow,0,1)$ be a residuated lattice. Moreover, let $f$ and $g$ be endomorphisms of $\mathcal{L}$.
An \emph{$(f,g)$-implicative derivation} of $\mathcal{L}$ is a function $d:L \rightarrow L$ which satisfies
\begin{align}
d(x \rightarrow y) = (d(x) \rightarrow f(y)) \vee (g(x) \rightarrow d(y))
\end{align}
for all $x,y \in L$.
\end{definition}

\section{Residuated Multilattices}
\label{remu0}
\begin{definition}[Residuated Multilattice]
A \emph{residuated multilattice} is an algebra $(M,\sqcup,\sqcap,\odot,\rightarrow,0,1)$ such that $(M,\sqcup,\sqcap,0,1)$ is a bounded multilattice,
$(M,\odot,1)$ is a commutative monoid, and $\rightarrow$ is a binary operation satisfying
\begin{align}
x \odot z \leq y \Leftrightarrow x \leq z \rightarrow y
\end{align}
for all $x,y,z \in M$.
\end{definition}

\begin{definition}[Homomorphism]
A \emph{homomorphism} of $f:(M_1,\sqcup,\sqcap,\odot,\rightarrow,0,1) \rightarrow (M_2,\sqcup,\sqcap,\odot,\rightarrow,0,1)$ of residuated multilattices
is a function $f:M_1 \rightarrow M_2$ which satisfies
\begin{align}
f(x \sqcap y) &\subseteq f(x) \sqcap f(y) \\
f(x \sqcup y) &\subseteq f(x) \sqcup f(y) \\
f(x \odot y) &= f(x) \odot f(y) \\
f(x \rightarrow y) &= f(x) \rightarrow f(y) \\
f(0) &= 0 \\
f(1) &= 1
\end{align}
for all $x,y \in M_1$.
\end{definition}

\begin{definition}[Multiplicative Derivation~\cite{YFL22}]
Let $\mathcal{M} := (M,\sqcup,\sqcap,\odot,\rightarrow,0,1)$ be a residuated multilattice. Moreover, let $f$ and $g$ be endomorphisms of $\mathcal{M}$.
An {$(f,g)$-multiplicative derivation} of $\mathcal{M}$ is a function $d:M \rightarrow M$ which satisfies
\begin{align}
d(x \odot y) \in (d(x) \odot f(y)) \sqcup (g(x) \odot d(y))
\end{align}
for all $x,y \in M$.
\end{definition}

\begin{definition}[Implicative Derivation]
Let $(M,\sqcup,\sqcap,\odot,\rightarrow,0,1)$ be a residuated multilattice. Moreover, let $f$ and $g$ be endomorphisms of $\mathcal{M}$.
An \emph{$(f,g)$-implicative derivation} of $\mathcal{M}$ is a function $d:M \rightarrow M$ which satisfies
\begin{align}
d(x \rightarrow y) \in (d(x) \rightarrow f(y)) \sqcup (g(x) \rightarrow d(y))
\end{align}
for all $x,y \in M$.
\end{definition}

\section*{Acknowledgements}
Thanks to Darline for feedback and error checking.

\bibliographystyle{plain}
\bibliography{specification}
\end{document}
